\documentclass[10pt,a4paper]{amsart}
\usepackage{amsmath,amssymb,amsthm,graphicx}
\usepackage{verbatim}
\newenvironment{reference}{\comment}{\endcomment}
%\newenvironment{reference}{}{}
\newenvironment{slogan}{\comment}{\endcomment}
\newenvironment{history}{\comment}{\endcomment}
% For commutative diagrams we use Xy-pic
\usepackage[all]{xy}

% We use 2cell for 2-commutative diagrams.
\xyoption{2cell}
% Theorem environments.
%
\theoremstyle{plain}
\newtheorem{theorem}[subsection]{Theorem}
\newtheorem{proposition}[subsection]{Proposition}
\newtheorem{lemma}[subsection]{Lemma}

\theoremstyle{definition}
\newtheorem{definition}[subsection]{Definition}
\newtheorem{example}[subsection]{Example}
\newtheorem{exercise}[subsection]{Exercise}
\newtheorem{situation}[subsection]{Situation}

\theoremstyle{remark}
\newtheorem{remark}[subsection]{Remark}
\newtheorem{remarks}[subsection]{Remarks}

\numberwithin{equation}{subsection}

% Macros
%
\def\lim{\mathop{\mathrm{lim}}\nolimits}
\def\colim{\mathop{\mathrm{colim}}\nolimits}
\def\Spec{\mathop{\mathrm{Spec}}}
\def\Hom{\mathop{\mathrm{Hom}}\nolimits}
\def\Ext{\mathop{\mathrm{Ext}}\nolimits}
\def\SheafHom{\mathop{\mathcal{H}\!\mathit{om}}\nolimits}
\def\SheafExt{\mathop{\mathcal{E}\!\mathit{xt}}\nolimits}
\def\Sch{\mathit{Sch}}
\def\Mor{\mathop{\mathrm{Mor}}\nolimits}
\def\Ob{\mathop{\mathrm{Ob}}\nolimits}
\def\Sh{\mathop{\mathit{Sh}}\nolimits}
\def\NL{\mathop{N\!L}\nolimits}
\def\CH{\mathop{\mathrm{CH}}\nolimits}
\def\proetale{{pro\text{-}\acute{e}tale}}
\def\etale{{\acute{e}tale}}
\def\QCoh{\mathit{QCoh}}
\def\Ker{\mathop{\mathrm{Ker}}}
\def\Im{\mathop{\mathrm{Im}}}
\def\Coker{\mathop{\mathrm{Coker}}}
\def\Coim{\mathop{\mathrm{Coim}}}

% Boxtimes
%
\DeclareMathSymbol{\boxtimes}{\mathbin}{AMSa}{"02}

%
% Macros for moduli stacks/spaces
%
\def\QCohstack{\mathcal{QC}\!\mathit{oh}}
\def\Cohstack{\mathcal{C}\!\mathit{oh}}
\def\Spacesstack{\mathcal{S}\!\mathit{paces}}
\def\Quotfunctor{\mathrm{Quot}}
\def\Hilbfunctor{\mathrm{Hilb}}
\def\Curvesstack{\mathcal{C}\!\mathit{urves}}
\def\Polarizedstack{\mathcal{P}\!\mathit{olarized}}
\def\Complexesstack{\mathcal{C}\!\mathit{omplexes}}
% \Pic is the operator that assigns to X its picard group, usage \Pic(X)
% \Picardstack_{X/B} denotes the Picard stack of X over B
% \Picardfunctor_{X/B} denotes the Picard functor of X over B
\def\Pic{\mathop{\mathrm{Pic}}\nolimits}
\def\Picardstack{\mathcal{P}\!\mathit{ic}}
\def\Picardfunctor{\mathrm{Pic}}
\def\Deformationcategory{\mathcal{D}\!\mathit{ef}}
\usepackage{cleveref}
\crefname{originalreference}{original reference}{original references}
\Crefname{originalreference}{Original reference}{Original references}
\makeatletter
\newlabel{algebraic-lemma-criterion-map-representable-spaces-fibred-in-groupoids}{{algebraic-lemma-criterion-map-representable-spaces-fibred-in-groupoids}{source}{}{sourcepointer-1}{}}
\newlabel{algebraic-lemma-criterion-map-representable-spaces-fibred-in-groupoids@cref}{{[originalreference][0][]algebraic-lemma-criterion-map-representable-spaces-fibred-in-groupoids}{[1][0][]source}{}{}{}}
\makeatother
\crefname{originalequation}{original equation}{original equations}
\Crefname{originalequation}{Original equation}{Original equations}
\makeatletter
\newlabel{equation-diagonal}{{equation-diagonal}{source}{}{sourcepointer-2}{}}
\newlabel{equation-diagonal@cref}{{[originalequation][0][]equation-diagonal}{[1][0][]source}{}{}{}}
\makeatother
\crefname{originallemma}{original lemma}{original lemmas}
\Crefname{originallemma}{Original lemma}{Original lemmas}
\makeatletter
\newlabel{lemma-properties-diagonal}{{lemma-properties-diagonal}{source}{}{sourcepointer-3}{}}
\newlabel{lemma-properties-diagonal@cref}{{[originallemma][0][]lemma-properties-diagonal}{[1][0][]source}{}{}{}}
\makeatother
\crefname{originalreference}{original reference}{original references}
\Crefname{originalreference}{Original reference}{Original references}
\makeatletter
\newlabel{spaces-morphisms-lemma-monomorphism-loc-finite-type-loc-quasi-finite}{{spaces-morphisms-lemma-monomorphism-loc-finite-type-loc-quasi-finite}{source}{}{sourcepointer-4}{}}
\newlabel{spaces-morphisms-lemma-monomorphism-loc-finite-type-loc-quasi-finite@cref}{{[originalreference][0][]spaces-morphisms-lemma-monomorphism-loc-finite-type-loc-quasi-finite}{[1][0][]source}{}{}{}}
\makeatother
\crefname{originalreference}{original reference}{original references}
\Crefname{originalreference}{Original reference}{Original references}
\makeatletter
\newlabel{spaces-morphisms-proposition-locally-quasi-finite-separated-over-scheme}{{spaces-morphisms-proposition-locally-quasi-finite-separated-over-scheme}{source}{}{sourcepointer-5}{}}
\newlabel{spaces-morphisms-proposition-locally-quasi-finite-separated-over-scheme@cref}{{[originalreference][0][]spaces-morphisms-proposition-locally-quasi-finite-separated-over-scheme}{[1][0][]source}{}{}{}}
\makeatother
\crefname{originalreference}{original reference}{original references}
\Crefname{originalreference}{Original reference}{Original references}
\makeatletter
\newlabel{stacks-properties-section-properties-morphisms}{{stacks-properties-section-properties-morphisms}{source}{}{sourcepointer-6}{}}
\newlabel{stacks-properties-section-properties-morphisms@cref}{{[originalreference][0][]stacks-properties-section-properties-morphisms}{[1][0][]source}{}{}{}}
\makeatother
\title{Three equivalent conditions for the diagonal of a morphism of algebraic stacks to be representable}
\author{The Stacks Project Authors}
\date{stacks-project:stacks-morphisms.tex:local-snapshot-20261012; GNU FDL 1.2 or later}
\begin{document}
\maketitle
\noindent\textit{Editorial note.} All delivered mathematical text is the author's own, copied from the pinned Stacks Project chapter source sources/stacks-algebraic-stacks-morphisms/pinned-original-source.tex; no statement or proof step was rewritten, reordered or removed.\par\medskip
\noindent\textit{Editorial note.} Preamble provenance: the Stacks Project keeps one single shared preamble for every chapter; this card takes it directly from the staged sources/stacks-algebraic-stacks-morphisms/preamble.tex as three verbatim spans (lines 8-12, 14-18 and 154-217), so no derived file was written and nothing was added or redefined. The wrapper class and the base packages are supplied by the assembler.\par\medskip
\noindent\textit{Editorial note.} Bibliography provenance: neither the counted statement nor its proof contains a citation (this lane rescans the whole span before writing the card), so the card carries no bibliography and no entry was reconstructed.\par\medskip
\noindent\textit{Editorial note.} Selection provenance: this card was opened under the file-based control contract (CONTROL.md), with the mechanical pre-card gate run on its statement and proof interval before the card was written.\par\medskip
\noindent\textit{Allowed standard prerequisites (editorial).}\begin{itemize}
\item The statement that the diagonal of any morphism of algebraic stacks is representable by algebraic spaces. Invoked at the start of the counted proof (source lines 419-422). Provides the ambient representability that makes the conditions comparable.
\item The criterion for a morphism of algebraic stacks representable by algebraic spaces to be a monomorphism. Used inside the counted proof (source lines 423-435). Links condition (2) with condition (3).
\item The description of the pullback of the diagonal along a test morphism as an isomorphism space. Used inside the counted proof (source lines 436-447). Turns the monomorphism condition into a condition on the fibres of the diagonal.
\end{itemize}\medskip
\ifdefined\hypertarget\hypertarget{sourcepointer-1}{}\fi
\noindent\textit{Source reference.} Original reference algebraic-lemma-criterion-map-representable-spaces-fibred-in-groupoids: sources/stacks-algebraic-stacks-morphisms/pinned-original-source.tex, source lines 431--431. Navigation pointer only: the counted text refers to the Stacks item algebraic-lemma-criterion-map-representable-spaces-fibred-in-groupoids, which belongs to another chapter of the project and is not part of this card. The pointer carries the target's own label as its number because this chapter's staged snapshot does not contain the other chapter's tag table; its content is not reproduced here.\par\medskip
\ifdefined\hypertarget\hypertarget{sourcepointer-2}{}\fi
\noindent\textit{Source reference.} Original equation equation-diagonal: sources/stacks-algebraic-stacks-morphisms/pinned-original-source.tex, source lines 375--375. Navigation pointer only: the counted text refers to the Stacks item equation-diagonal, which is an existing labeled result of this same chapter and is not part of this card. The pointer carries the target's own label as its number and its definition line as its span; its content is not reproduced here.\par\medskip
\ifdefined\hypertarget\hypertarget{sourcepointer-3}{}\fi
\noindent\textit{Source reference.} Original lemma lemma-properties-diagonal: sources/stacks-algebraic-stacks-morphisms/pinned-original-source.tex, source lines 344--344. Navigation pointer only: the counted text refers to the Stacks item lemma-properties-diagonal, which is an existing labeled result of this same chapter and is not part of this card. The pointer carries the target's own label as its number and its definition line as its span; its content is not reproduced here.\par\medskip
\ifdefined\hypertarget\hypertarget{sourcepointer-4}{}\fi
\noindent\textit{Source reference.} Original reference spaces-morphisms-lemma-monomorphism-loc-finite-type-loc-quasi-finite: sources/stacks-algebraic-stacks-morphisms/pinned-original-source.tex, source lines 439--439. Navigation pointer only: the counted text refers to the Stacks item spaces-morphisms-lemma-monomorphism-loc-finite-type-loc-quasi-finite, which belongs to another chapter of the project and is not part of this card. The pointer carries the target's own label as its number because this chapter's staged snapshot does not contain the other chapter's tag table; its content is not reproduced here.\par\medskip
\ifdefined\hypertarget\hypertarget{sourcepointer-5}{}\fi
\noindent\textit{Source reference.} Original reference spaces-morphisms-proposition-locally-quasi-finite-separated-over-scheme: sources/stacks-algebraic-stacks-morphisms/pinned-original-source.tex, source lines 445--445. Navigation pointer only: the counted text refers to the Stacks item spaces-morphisms-proposition-locally-quasi-finite-separated-over-scheme, which belongs to another chapter of the project and is not part of this card. The pointer carries the target's own label as its number because this chapter's staged snapshot does not contain the other chapter's tag table; its content is not reproduced here.\par\medskip
\ifdefined\hypertarget\hypertarget{sourcepointer-6}{}\fi
\noindent\textit{Source reference.} Original reference stacks-properties-section-properties-morphisms: sources/stacks-algebraic-stacks-morphisms/pinned-original-source.tex, source lines 248--248. Navigation pointer only: the counted text refers to the Stacks item stacks-properties-section-properties-morphisms, which belongs to another chapter of the project and is not part of this card. The pointer carries the target's own label as its number because this chapter's staged snapshot does not contain the other chapter's tag table; its content is not reproduced here.\par\medskip
\begin{lemma}
\label{lemma-properties-diagonal-representable}
Let $f : \mathcal{X} \to \mathcal{Y}$ be a morphism of algebraic stacks
which is representable by algebraic spaces. Then
\begin{enumerate}
\item $\Delta_f$ is representable
(by schemes),
\item $\Delta_f$ is locally of finite type,
\item $\Delta_f$ is a monomorphism,
\item $\Delta_f$ is separated, and
\item $\Delta_f$ is locally quasi-finite.
\end{enumerate}
\end{lemma}

\begin{proof}
We have already seen in
Lemma \ref{lemma-properties-diagonal}
that $\Delta_f$ is representable by algebraic
spaces. Hence the statements (2) -- (5) make sense, see
Properties of Stacks,
Section \ref{stacks-properties-section-properties-morphisms}.
Also
Lemma \ref{lemma-properties-diagonal}
guarantees (2) holds.
Let $T \to \mathcal{X} \times_\mathcal{Y} \mathcal{X}$ be a morphism
and contemplate Diagram (\ref{equation-diagonal}). By
Algebraic Stacks, Lemma
\ref{algebraic-lemma-criterion-map-representable-spaces-fibred-in-groupoids}
the right vertical arrow is injective as a map of sheaves, i.e., a
monomorphism of algebraic spaces. Hence also the morphism
$T \times_{\mathcal{X} \times_\mathcal{Y} \mathcal{X}} \mathcal{X} \to T$
is a monomorphism. Thus (3) holds. We already know that
$T \times_{\mathcal{X} \times_\mathcal{Y} \mathcal{X}} \mathcal{X} \to T$
is locally of finite type. Thus
Morphisms of Spaces, Lemma
\ref{spaces-morphisms-lemma-monomorphism-loc-finite-type-loc-quasi-finite}
allows us to conclude that
$T \times_{\mathcal{X} \times_\mathcal{Y} \mathcal{X}} \mathcal{X} \to T$
is locally quasi-finite and separated. This proves (4) and (5).
Finally,
Morphisms of Spaces, Proposition
\ref{spaces-morphisms-proposition-locally-quasi-finite-separated-over-scheme}
implies that
$T \times_{\mathcal{X} \times_\mathcal{Y} \mathcal{X}} \mathcal{X}$
is a scheme which proves (1).
\end{proof}

\end{document}
